\documentclass[10pt,a4paper]{amsart}
\usepackage[margin=19mm]{geometry}
\usepackage{amsmath,amssymb,mathtools}
\usepackage[scr=rsfs,cal=euler]{mathalfa}
\usepackage{xfrac}
\usepackage[unicode,hidelinks,bookmarks=false]{hyperref}
\newcommand{\ie}{i.e.\ }
\newcommand{\cf}{cf.\ }
\newcommand{\resp}{respectively\ }
\newcommand{\Subsection}{\subsection}
\providecommand{\nobreakdash}{\nobreak\hbox{-}}
\setlength{\emergencystretch}{2em}
\multlinegap0pt
\usepackage{amssymb}
\usepackage{dsfont}
\newtheorem{lemma}{Lemma}
\newtheorem{theorem}{Theorem}
\title{On a slight weakening of Kripke-Platek Set Theory}
\author{Zachiri McKenzie}
\date{Source: arXiv:2608.23398v1 (CC BY 4.0)}
\begin{document}
\maketitle

\noindent\emph{Source and licence.} On a slight weakening of Kripke-Platek Set Theory. Version arXiv:2608.23398v1 (CC BY 4.0), distributed by arXiv under the Creative Commons Attribution 4.0 International licence, http://creativecommons.org/licenses/by/4.0/, as recorded in the arXiv licence metadata for this exact version and shown on the abstract page https://arxiv.org/abs/2608.23398v1. Full licence notice: \texttt{licenses/arxiv-2608.23398-CC-BY-4.0.txt}.\par\smallskip

\noindent\textit{Editorial scope.} Counted unit: the article's theorem theorem (source0.tex lines 694--696). The article's complete original proof of it is delivered with it (source0.tex lines 698--714, one \texttt{proof} environment of 17 lines). No mathematics is written, repaired or re-derived: the statement and the proof are the article's own lines, in the article's own notation, and no retained line is substituted or re-flowed.  Retained uncounted as the source-local context the proof needs: lines 640--642; lines 658--660; lines 666--668; lines 678--680.  Nothing was omitted from any retained span, so the package carries no omission record.  No retained line contains a citation, so the wrapper carries no bibliography and no bibliography entry.  No retained line contains a figure environment, an image-inclusion command or drawing code, so no image is delivered and the package declares no asset dependency.  Editorial formatting only, in addition: an AMS \texttt{amsart} preamble that restates the article's own declarations and macros used by the retained lines, namely the environments \texttt{lemma}, \texttt{theorem} on separate counters, together with the standard packages those lines need.  The retained environments print in this document as Lemma 1, Theorem 1 in order of appearance, whereas the article prints the same items with the numbering of its own full document.  The complete native source of the article as distributed in the arXiv e-print for this exact version is bundled unchanged as \texttt{sources/208408/source0.tex}.
\begin{lemma} \label{Th:KSuperTransitiveSubstructure}
$\mathcal{K} \subseteq_{\mathsf{tr}}^{\mathcal{P}} \mathcal{M}$.
\end{lemma} 

\begin{lemma} \label{Th:PowersetsInK}
For all $x \in K$, $\mathcal{P}^{\mathcal{M}}(x) \in K$.
\end{lemma}

\begin{lemma} \label{Th:UnionsInK}
For all $x \in K$, $\bigcup x \in K$.
\end{lemma}

\begin{lemma} \label{Th:PairingInK}
For all $x, y \in K$, $\{x, y\} \in K$.
\end{lemma}

\begin{theorem}
The structure $\mathcal{K}$ satisfies $\mathsf{S}_0+\textsf{Infinity}+\textsf{Powerset}+\Delta_0^{\mathcal{P}}\textsf{-Separation}+\Delta_0^{\mathcal{P}}\textsf{-Collection}+\textsf{Set-Foundation}+\mathsf{AC}$.
\end{theorem}

\begin{proof}
Lemmas \ref{Th:KSuperTransitiveSubstructure}, \ref{Th:PowersetsInK}, \ref{Th:UnionsInK} and \ref{Th:PairingInK} immediately imply that $\mathcal{K}$ satisfies $\mathsf{S}_0+\textsf{Powerset}+\Delta_0^{\mathcal{P}}\textsf{-Separation}+\textsf{Set-Foundation}+\mathsf{AC}$. Since $\mathcal{M} \models (\omega \subseteq K_0)$, $\mathcal{M} \models (\omega \in K_1)$ and $\omega^{\mathcal{M}} \in K$. Therefore, $\textsf{Infinity}$ holds in $\mathcal{K}$. We are left to verify that $\Delta_0^{\mathcal{P}}\textsf{-Collection}$ holds in $\mathcal{K}$. Towards this end, let $\phi(x, y, \vec{z})$ be a $\Delta_0^{\mathcal{P}}$-formula. Let $b, \vec{a} \in K$ be such that 
\begin{equation} \label{eq:CollectionAntecedent}
\mathcal{K} \models (\forall x \in b)\exists y \phi(x, y, \vec{a}).
\end{equation}
Work inside $\mathcal{M}$. Let
\[
q= \{ \langle x, n \rangle \in b \times \omega \mid (\exists y \in K_n) \phi(x, y, \vec{a}) \land (\forall m \in n)(\forall z \in K_m) \neg \phi(x, z, \vec{a})\},
\]
and let $d= \mathsf{rng}(q)$.

Work in the metatheory again. Now, by (\ref{eq:CollectionAntecedent}), $d^* \subseteq \omega$. Therefore, there exists $l \in \omega$ such that for all $n \in d^*$, $n < l$, otherwise the standard natural numbers would be a set in $\mathcal{M}$. So, 
\[
\mathcal{K} \models (\forall x \in b)(\exists y \in K_l)\phi(x, y, \vec{a}).
\]
This shows that $\Delta_0^{\mathcal{P}}\textsf{-Collection}$ holds in $\mathcal{K}$. \hfill$\square$  
\end{proof}

\end{document}
